\documentclass[12pt,a4paper]{article}

% Packages
\usepackage[utf8]{inputenc}
\usepackage[english]{babel}
\usepackage{amsmath,amsthm,amssymb,amsfonts}
\usepackage{mathtools}
\usepackage{physics}
\usepackage{geometry}
\usepackage{fancyhdr}
\usepackage{graphicx}
\usepackage{hyperref}
\usepackage{tcolorbox}
\usepackage{xcolor}
\usepackage{enumitem}
\usepackage{tikz}
\usepackage{pgfplots}
\pgfplotsset{compat=1.17}

% Page geometry
\geometry{margin=1in}

% Header and footer
\pagestyle{fancy}
\fancyhf{}
\lhead{GATE 2026 Classical Mechanics Study Guide}
\rhead{\thepage}
\renewcommand{\headrulewidth}{0.5pt}

% Custom theorem environments
\theoremstyle{definition}
\newtheorem{theorem}{Theorem}[section]
\newtheorem{definition}[theorem]{Definition}
\newtheorem{lemma}[theorem]{Lemma}
\newtheorem{corollary}[theorem]{Corollary}
\newtheorem{example}[theorem]{Example}
\newtheorem{formula}[theorem]{Formula}

% Custom boxes
\tcbuselibrary{theorems}
\newtcolorbox{importantbox}{
    colback=red!5!white,
    colframe=red!75!black,
    title=Important for GATE
}

\newtcolorbox{tipbox}{
    colback=blue!5!white,
    colframe=blue!75!black,
    title=Exam Tip
}

\newtcolorbox{notebox}{
    colback=yellow!5!white,
    colframe=orange!75!black,
    title=Note
}

% Title page
\title{\Huge\textbf{Classical Mechanics}\\[0.3cm]
\Large Complete Study Guide for GATE 2026\\[0.2cm]
\large Problem-Solving Intensive Mastery}
\author{Comprehensive Coverage of All Topics}
\date{Prepared: December 2025}

\begin{document}

\maketitle
\newpage

\tableofcontents
\newpage

%======================================================================
% INTRODUCTION
%======================================================================
\section{Introduction}

\begin{importantbox}
Classical Mechanics is the most problem-solving intensive subject in GATE Physics. It requires not just understanding but the ability to apply concepts quickly and accurately. This guide covers every essential topic with complete problem-solving strategies and NO STONE LEFT UNTURNED.
\end{importantbox}

\subsection{Why Classical Mechanics is Crucial}
\begin{enumerate}
    \item Foundation for Quantum Mechanics and Statistical Mechanics
    \item Most predictable question patterns in GATE
    \item High weightage (15-20\% of physics questions)
    \item Directly applicable problem-solving skills
    \item Tests mathematical physics mastery
\end{enumerate}

\subsection{Study Strategy}
\begin{tipbox}
\begin{itemize}
    \item Master the formulation before solving problems
    \item Practice at least 50 problems per topic
    \item Time yourself - aim for 3-4 minutes per problem
    \item Understand physical meaning, not just equations
    \item Connect Lagrangian and Hamiltonian viewpoints
\end{itemize}
\end{tipbox}

%======================================================================
% CHAPTER 1: LAGRANGIAN FORMULATION
%======================================================================
\newpage
\section{Lagrangian Formulation}

\subsection{Motivation and Fundamental Concepts}

\begin{definition}
The \textbf{Lagrangian} of a system is defined as:
\[
L = T - V
\]
where $T$ is kinetic energy and $V$ is potential energy.
\end{definition}

\begin{importantbox}
\textbf{Why Lagrangian Mechanics?}
\begin{itemize}
    \item Works with any set of coordinates (not just Cartesian)
    \item Constraints handled automatically
    \item Scalar function (easier than vector equations)
    \item Suitable for complex systems
    \item Basis for quantum field theory
\end{itemize}
\end{importantbox}

\subsection{Generalized Coordinates}

\begin{definition}
\textbf{Generalized coordinates} $q_1, q_2, \ldots, q_n$ are any set of independent coordinates that completely specify the configuration of a system.

\textbf{Degrees of freedom:} Number of independent generalized coordinates needed.

For $N$ particles in 3D with $k$ holonomic constraints:
\[
\text{DOF} = 3N - k
\]
\end{definition}

\begin{example}
Find degrees of freedom for:
\begin{enumerate}
    \item Single particle in space: DOF = 3 (can use $x, y, z$ or $r, \theta, \phi$)
    \item Particle on sphere: DOF = 2 (constraint: $x^2 + y^2 + z^2 = R^2$)
    \item Double pendulum: DOF = 2 (use $\theta_1, \theta_2$)
    \item Rigid body: DOF = 6 (3 translation + 3 rotation)
\end{enumerate}
\end{example}

\subsection{D'Alembert's Principle and Virtual Work}

\begin{definition}
\textbf{Virtual displacement} $\delta q_i$: Infinitesimal displacement consistent with constraints at fixed time.

\textbf{Virtual work:}
\[
\delta W = \sum_i F_i \cdot \delta r_i
\]
\end{definition}

\begin{theorem}[D'Alembert's Principle]
For a system in equilibrium under applied forces $F_i$ and constraint forces:
\[
\sum_i (F_i - \dot{p}_i) \cdot \delta r_i = 0
\]
where constraint forces do no virtual work.
\end{theorem}

\subsection{Euler-Lagrange Equations}

\begin{theorem}[Euler-Lagrange Equations]
For a system with Lagrangian $L(q_i, \dot{q}_i, t)$, the equations of motion are:
\[
\frac{d}{dt}\left(\frac{\partial L}{\partial \dot{q}_i}\right) - \frac{\partial L}{\partial q_i} = 0
\]
for each generalized coordinate $q_i$.
\end{theorem}

\begin{importantbox}
\textbf{Steps to Solve Problems Using Lagrangian:}
\begin{enumerate}
    \item Choose appropriate generalized coordinates
    \item Write kinetic energy $T$ in terms of $q_i$ and $\dot{q}_i$
    \item Write potential energy $V$ in terms of $q_i$
    \item Form Lagrangian: $L = T - V$
    \item Apply Euler-Lagrange equation for each coordinate
    \item Solve the resulting differential equations
\end{enumerate}
\end{importantbox}

\begin{example}
\textbf{Simple Pendulum using Lagrangian Formulation}

Setup: Mass $m$, length $l$, angle $\theta$ from vertical.

\textbf{Solution:}
\begin{enumerate}
    \item Generalized coordinate: $q = \theta$
    
    \item Kinetic energy:
    \[
    T = \frac{1}{2}m v^2 = \frac{1}{2}m(l\dot{\theta})^2 = \frac{1}{2}ml^2\dot{\theta}^2
    \]
    
    \item Potential energy (taking lowest point as reference):
    \[
    V = mgl(1 - \cos\theta)
    \]
    
    \item Lagrangian:
    \[
    L = T - V = \frac{1}{2}ml^2\dot{\theta}^2 - mgl(1 - \cos\theta)
    \]
    
    \item Apply Euler-Lagrange:
    \[
    \frac{\partial L}{\partial \theta} = -mgl\sin\theta, \quad \frac{\partial L}{\partial \dot{\theta}} = ml^2\dot{\theta}
    \]
    \[
    \frac{d}{dt}(ml^2\dot{\theta}) = -mgl\sin\theta
    \]
    
    \item Equation of motion:
    \[
    \ddot{\theta} + \frac{g}{l}\sin\theta = 0
    \]
    
    For small angles: $\ddot{\theta} + \omega_0^2\theta = 0$ where $\omega_0 = \sqrt{g/l}$
\end{enumerate}
\end{example}

\subsection{Constraints}

\begin{definition}
\textbf{Holonomic constraints:} Can be written as $f(q_1, \ldots, q_n, t) = 0$

Examples:
\begin{itemize}
    \item Particle on sphere: $x^2 + y^2 + z^2 - R^2 = 0$
    \item Rigid body: $|r_i - r_j| = \text{const}$
    \item Rolling without slipping: $x = R\theta$
\end{itemize}

\textbf{Non-holonomic constraints:} Involve velocities, cannot be integrated.

Example: Rolling without slipping in general: $v_{\text{contact}} = 0$
\end{definition}

\begin{tipbox}
\textbf{Handling Constraints:}
\begin{itemize}
    \item \textbf{Holonomic:} Reduce number of coordinates OR use Lagrange multipliers
    \item \textbf{Non-holonomic:} Must use Lagrange multipliers
\end{itemize}
\end{tipbox}

\subsection{Lagrange Multipliers Method}

\begin{theorem}
For constraints $g_k(q, t) = 0$, the modified Euler-Lagrange equations are:
\[
\frac{d}{dt}\left(\frac{\partial L}{\partial \dot{q}_i}\right) - \frac{\partial L}{\partial q_i} = \sum_k \lambda_k \frac{\partial g_k}{\partial q_i}
\]
where $\lambda_k$ are Lagrange multipliers (related to constraint forces).
\end{theorem}

\begin{example}
\textbf{Bead on Rotating Hoop}

A bead of mass $m$ slides on a hoop of radius $R$ rotating with constant angular velocity $\omega$ about vertical diameter.

\textbf{Solution:}
Generalized coordinate: $\theta$ (angle from top)

Position: $x = R\sin\theta\cos(\omega t)$, $y = R\sin\theta\sin(\omega t)$, $z = R\cos\theta$

Velocities:
\begin{align*}
\dot{x} &= R\dot{\theta}\cos\theta\cos(\omega t) - R\omega\sin\theta\sin(\omega t) \\
\dot{y} &= R\dot{\theta}\cos\theta\sin(\omega t) + R\omega\sin\theta\cos(\omega t) \\
\dot{z} &= -R\dot{\theta}\sin\theta
\end{align*}

Kinetic energy:
\begin{align*}
T &= \frac{1}{2}m(\dot{x}^2 + \dot{y}^2 + \dot{z}^2) \\
&= \frac{1}{2}m[R^2\dot{\theta}^2 + R^2\omega^2\sin^2\theta]
\end{align*}

Potential energy: $V = -mgR\cos\theta$

Lagrangian:
\[
L = \frac{1}{2}mR^2[\dot{\theta}^2 + \omega^2\sin^2\theta] + mgR\cos\theta
\]

Euler-Lagrange equation:
\[
\frac{d}{dt}(mR^2\dot{\theta}) - mR^2\omega^2\sin\theta\cos\theta - mgR\sin\theta = 0
\]

\[
\ddot{\theta} = (\omega^2\cos\theta - \frac{g}{R})\sin\theta
\]

Equilibrium: $\ddot{\theta} = 0 \implies \cos\theta = \frac{g}{R\omega^2}$ (for $\omega^2 > g/R$)
\end{example}

\subsection{Generalized Momentum and Cyclic Coordinates}

\begin{definition}
\textbf{Generalized momentum} (canonical momentum):
\[
p_i = \frac{\partial L}{\partial \dot{q}_i}
\]

\textbf{Cyclic (ignorable) coordinate:} If $\frac{\partial L}{\partial q_i} = 0$, then $q_i$ is cyclic.
\end{definition}

\begin{theorem}
If $q_i$ is a cyclic coordinate, then its conjugate momentum $p_i$ is conserved:
\[
\frac{dp_i}{dt} = 0 \implies p_i = \text{constant}
\]
\end{theorem}

\begin{importantbox}
\textbf{Common Cyclic Coordinates and Conservation Laws:}
\begin{itemize}
    \item Translation invariance ($x$ cyclic) $\implies$ Linear momentum conserved
    \item Rotational invariance ($\theta$ cyclic) $\implies$ Angular momentum conserved
    \item Time independence ($t$ not explicit) $\implies$ Energy conserved
\end{itemize}
\end{importantbox}

\begin{example}
\textbf{Central Force Problem in Lagrangian Formulation}

For a particle in central force $F(r)$ (potential $V(r)$):

Using polar coordinates $(r, \theta)$:
\[
L = \frac{1}{2}m(\dot{r}^2 + r^2\dot{\theta}^2) - V(r)
\]

$\theta$ is cyclic! Therefore:
\[
p_\theta = \frac{\partial L}{\partial \dot{\theta}} = mr^2\dot{\theta} = l = \text{const (angular momentum)}
\]

For $r$:
\[
\frac{d}{dt}(m\dot{r}) - mr\dot{\theta}^2 + \frac{dV}{dr} = 0
\]

Using $l = mr^2\dot{\theta}$:
\[
m\ddot{r} = \frac{l^2}{mr^3} - \frac{dV}{dr}
\]
\end{example}

\subsection{Conservation Laws and Noether's Theorem}

\begin{theorem}[Noether's Theorem]
Every continuous symmetry of the Lagrangian corresponds to a conservation law:
\begin{itemize}
    \item \textbf{Translation symmetry} $\implies$ Momentum conservation
    \item \textbf{Rotation symmetry} $\implies$ Angular momentum conservation
    \item \textbf{Time translation} $\implies$ Energy conservation
\end{itemize}
\end{theorem}

\subsection{Energy Function and Conservation}

\begin{definition}
The \textbf{energy function} (Jacobi's integral):
\[
h = \sum_i \dot{q}_i \frac{\partial L}{\partial \dot{q}_i} - L = \sum_i p_i\dot{q}_i - L
\]
\end{definition}

\begin{theorem}
If $L$ does not depend explicitly on time ($\partial L/\partial t = 0$), then energy is conserved:
\[
\frac{dh}{dt} = -\frac{\partial L}{\partial t}
\]

If $\partial L/\partial t = 0$, then $h = E = T + V$ (total energy).
\end{theorem}

\subsection{Practice Problems - Lagrangian Mechanics}

\begin{enumerate}
    \item Derive equations of motion for Atwood machine using Lagrangian.
    
    \item Find Lagrangian for double pendulum and derive equations of motion.
    
    \item A cylinder rolls without slipping down an incline. Use Lagrangian with constraint.
    
    \item Particle on surface of sphere under gravity - find equation of motion.
    
    \item Two masses connected by spring on frictionless table - find normal modes.
\end{enumerate}

%======================================================================
% CHAPTER 2: HAMILTONIAN MECHANICS
%======================================================================
\newpage
\section{Hamiltonian Mechanics}

\subsection{From Lagrangian to Hamiltonian}

\begin{definition}
The \textbf{Hamiltonian} is obtained via Legendre transformation:
\[
H(q_i, p_i, t) = \sum_i p_i\dot{q}_i - L(q_i, \dot{q}_i, t)
\]
where $p_i = \frac{\partial L}{\partial \dot{q}_i}$ are generalized momenta.
\end{definition}

\begin{importantbox}
\textbf{Procedure to Find Hamiltonian:}
\begin{enumerate}
    \item Start with Lagrangian $L(q, \dot{q}, t)$
    \item Find generalized momenta: $p_i = \frac{\partial L}{\partial \dot{q}_i}$
    \item Invert to get $\dot{q}_i = \dot{q}_i(q, p, t)$
    \item Substitute into $H = \sum p_i\dot{q}_i - L$
    \item Express everything in terms of $(q, p, t)$
\end{enumerate}
\end{importantbox}

\begin{example}
\textbf{Hamiltonian for Simple Harmonic Oscillator}

Lagrangian: $L = \frac{1}{2}m\dot{x}^2 - \frac{1}{2}kx^2$

Momentum: $p = \frac{\partial L}{\partial \dot{x}} = m\dot{x} \implies \dot{x} = \frac{p}{m}$

Hamiltonian:
\begin{align*}
H &= p\dot{x} - L = p \cdot \frac{p}{m} - \frac{1}{2}m\left(\frac{p}{m}\right)^2 + \frac{1}{2}kx^2 \\
&= \frac{p^2}{m} - \frac{p^2}{2m} + \frac{1}{2}kx^2 \\
&= \frac{p^2}{2m} + \frac{1}{2}kx^2 = T + V
\end{align*}
\end{example}

\subsection{Hamilton's Equations}

\begin{theorem}[Hamilton's Canonical Equations]
\[
\dot{q}_i = \frac{\partial H}{\partial p_i}, \quad \dot{p}_i = -\frac{\partial H}{\partial q_i}
\]
These are $2n$ first-order equations (versus $n$ second-order Euler-Lagrange equations).
\end{theorem}

\begin{tipbox}
\textbf{When to Use Hamiltonian vs Lagrangian:}
\begin{itemize}
    \item \textbf{Lagrangian:} Better for deriving equations of motion
    \item \textbf{Hamiltonian:} Better for phase space analysis, perturbation theory, quantum mechanics transition
    \item Both give same physics!
\end{itemize}
\end{tipbox}

\begin{example}
\textbf{Hamilton's Equations for SHO}

From previous: $H = \frac{p^2}{2m} + \frac{1}{2}kx^2$

\[
\dot{x} = \frac{\partial H}{\partial p} = \frac{p}{m}
\]
\[
\dot{p} = -\frac{\partial H}{\partial x} = -kx
\]

Combining: $m\ddot{x} = \dot{p} = -kx \implies \ddot{x} + \omega^2 x = 0$ where $\omega = \sqrt{k/m}$
\end{example}

\subsection{Phase Space}

\begin{definition}
\textbf{Phase space} is the $2n$-dimensional space of all $(q_i, p_i)$ coordinates.

\textbf{Phase space trajectory:} Path traced by system in phase space as it evolves.

\textbf{Phase portrait:} Collection of trajectories for different initial conditions.
\end{definition}

\begin{importantbox}
\textbf{Properties of Phase Space Trajectories:}
\begin{enumerate}
    \item Trajectories never cross (deterministic evolution)
    \item Closed trajectory $\implies$ Periodic motion
    \item Energy surfaces are constant $H$ surfaces
    \item Liouville's theorem: Phase space volume conserved
\end{enumerate}
\end{importantbox}

\subsection{Poisson Brackets}

\begin{definition}
The \textbf{Poisson bracket} of two functions $f(q, p, t)$ and $g(q, p, t)$ is:
\[
\{f, g\} = \sum_i \left(\frac{\partial f}{\partial q_i}\frac{\partial g}{\partial p_i} - \frac{\partial f}{\partial p_i}\frac{\partial g}{\partial q_i}\right)
\]
\end{definition}

\begin{importantbox}
\textbf{Fundamental Poisson Brackets (FPB):}
\[
\{q_i, q_j\} = 0, \quad \{p_i, p_j\} = 0, \quad \{q_i, p_j\} = \delta_{ij}
\]
\end{importantbox}

\begin{theorem}[Time Evolution via Poisson Brackets]
For any function $f(q, p, t)$:
\[
\frac{df}{dt} = \{f, H\} + \frac{\partial f}{\partial t}
\]

In particular, Hamilton's equations are:
\[
\dot{q}_i = \{q_i, H\}, \quad \dot{p}_i = \{p_i, H\}
\]
\end{theorem}

\begin{theorem}[Conservation and Poisson Brackets]
A quantity $f$ is conserved if:
\[
\frac{df}{dt} = 0 \iff \{f, H\} = 0 \text{ and } \frac{\partial f}{\partial t} = 0
\]
\end{theorem}

\begin{example}
\textbf{Angular Momentum Conservation}

For central force, $H = \frac{p^2}{2m} + V(r)$ where $r = \sqrt{x^2 + y^2 + z^2}$

Angular momentum: $L_z = xp_y - yp_x$

Check conservation:
\begin{align*}
\{L_z, H\} &= \{xp_y - yp_x, \frac{p_x^2 + p_y^2 + p_z^2}{2m} + V(r)\} \\
&= \{xp_y, V(r)\} - \{yp_x, V(r)\} \\
&= x\frac{\partial V}{\partial y} - y\frac{\partial V}{\partial x} \\
&= x \cdot \frac{dV}{dr} \cdot \frac{y}{r} - y \cdot \frac{dV}{dr} \cdot \frac{x}{r} = 0
\end{align*}

Therefore $L_z$ is conserved!
\end{example}

\subsection{Properties of Poisson Brackets}

\begin{formula}[Poisson Bracket Properties]
\begin{enumerate}
    \item \textbf{Antisymmetry:} $\{f, g\} = -\{g, f\}$
    
    \item \textbf{Linearity:} $\{af + bg, h\} = a\{f, h\} + b\{g, h\}$
    
    \item \textbf{Leibniz rule:} $\{fg, h\} = f\{g, h\} + \{f, h\}g$
    
    \item \textbf{Jacobi identity:} $\{f, \{g, h\}\} + \{g, \{h, f\}\} + \{h, \{f, g\}\} = 0$
\end{enumerate}
\end{formula}

\subsection{Canonical Transformations}

\begin{definition}
A transformation $(q, p) \to (Q, P)$ is \textbf{canonical} if it preserves the form of Hamilton's equations:
\[
\dot{Q}_i = \frac{\partial K}{\partial P_i}, \quad \dot{P}_i = -\frac{\partial K}{\partial Q_i}
\]
where $K(Q, P, t)$ is the new Hamiltonian.
\end{definition}

\begin{theorem}[Conditions for Canonical Transformation]
A transformation is canonical if and only if:
\begin{enumerate}
    \item It preserves Poisson brackets (symplectic condition):
    \[
    \{Q_i, Q_j\} = 0, \quad \{P_i, P_j\} = 0, \quad \{Q_i, P_j\} = \delta_{ij}
    \]
    
    \item OR there exists a generating function $F$ such that:
    \[
    p_i dq_i - P_i dQ_i = dF
    \]
\end{enumerate}
\end{theorem}

\subsection{Generating Functions}

\begin{importantbox}
\textbf{Four Types of Generating Functions:}

\textbf{Type 1: $F_1(q, Q, t)$}
\[
p_i = \frac{\partial F_1}{\partial q_i}, \quad P_i = -\frac{\partial F_1}{\partial Q_i}, \quad K = H + \frac{\partial F_1}{\partial t}
\]

\textbf{Type 2: $F_2(q, P, t)$}
\[
p_i = \frac{\partial F_2}{\partial q_i}, \quad Q_i = \frac{\partial F_2}{\partial P_i}, \quad K = H + \frac{\partial F_2}{\partial t}
\]

\textbf{Type 3: $F_3(p, Q, t)$}
\[
q_i = -\frac{\partial F_3}{\partial p_i}, \quad P_i = -\frac{\partial F_3}{\partial Q_i}, \quad K = H + \frac{\partial F_3}{\partial t}
\]

\textbf{Type 4: $F_4(p, P, t)$}
\[
q_i = -\frac{\partial F_4}{\partial p_i}, \quad Q_i = \frac{\partial F_4}{\partial P_i}, \quad K = H + \frac{\partial F_4}{\partial t}
\]
\end{importantbox}

\begin{example}
\textbf{Identity Transformation}

$F_2 = q_i P_i$

\[
p_i = \frac{\partial F_2}{\partial q_i} = P_i, \quad Q_i = \frac{\partial F_2}{\partial P_i} = q_i
\]

This swaps coordinates and momenta!
\end{example}

\begin{example}
\textbf{Point Transformation}

$F_2 = f(q, t) P$ generates $Q = f(q, t)$, $P = p \frac{\partial q}{\partial Q}$
\end{example}

\subsection{Hamilton-Jacobi Theory}

\begin{definition}
The \textbf{Hamilton-Jacobi equation} is:
\[
H\left(q_i, \frac{\partial S}{\partial q_i}, t\right) + \frac{\partial S}{\partial t} = 0
\]
where $S(q, t)$ is Hamilton's principal function.
\end{definition}

\begin{theorem}
If $S(q, \alpha, t)$ is a complete solution with $n$ independent constants $\alpha_i$, then:
\[
p_i = \frac{\partial S}{\partial q_i}, \quad \beta_i = \frac{\partial S}{\partial \alpha_i}
\]
give the solution with $\alpha_i$ and $\beta_i$ as constants of integration.
\end{theorem}

\begin{tipbox}
\textbf{Hamilton-Jacobi Method:}
\begin{enumerate}
    \item Write Hamilton-Jacobi equation
    \item Try separation: $S = W(q) - Et$ for time-independent $H$
    \item Solve for $W(q)$: $H(q, \partial W/\partial q) = E$
    \item Find $p = \partial S/\partial q$ and integrate for $q(t)$
\end{enumerate}
\end{tipbox}

\begin{example}
\textbf{Harmonic Oscillator via Hamilton-Jacobi}

$H = \frac{p^2}{2m} + \frac{1}{2}kx^2 = E$

Try $S = W(x) - Et$:
\[
\frac{1}{2m}\left(\frac{dW}{dx}\right)^2 + \frac{1}{2}kx^2 = E
\]

\[
\frac{dW}{dx} = \sqrt{2m(E - \frac{1}{2}kx^2)}
\]

\[
W = \int \sqrt{2mE - mkx^2}\, dx
\]

Complete solution leads to $x(t) = A\cos(\omega t + \phi)$ where $\omega = \sqrt{k/m}$
\end{example}

\subsection{Action-Angle Variables}

\begin{definition}
For periodic motion, \textbf{action variable}:
\[
J = \oint p\, dq = \oint p(q, E)\, dq
\]
where integral is over one complete period.

\textbf{Angle variable} $w$ conjugate to $J$ increases by $2\pi$ per cycle.
\end{definition}

\begin{theorem}
In action-angle variables:
\[
H = H(J) \implies \dot{J} = 0, \quad \dot{w} = \frac{\partial H}{\partial J} = \omega(J)
\]
Motion is uniform in $w$: $w = \omega t + w_0$
\end{theorem}

\subsection{Practice Problems - Hamiltonian Mechanics}

\begin{enumerate}
    \item Find Hamiltonian for particle in electromagnetic field.
    
    \item Calculate $\{L_x, L_y\}$ where $\vec{L} = \vec{r} \times \vec{p}$.
    
    \item Show that $F_2 = qP\cot\theta$ generates canonical transformation.
    
    \item Solve simple pendulum using Hamilton-Jacobi equation.
    
    \item Find action-angle variables for 1D harmonic oscillator.
\end{enumerate}

%======================================================================
% CHAPTER 3: CENTRAL FORCE MOTION
%======================================================================
\newpage
\section{Central Force Motion}

\subsection{Central Force Definition}

\begin{definition}
A \textbf{central force} is one that:
\begin{itemize}
    \item Depends only on distance from origin: $\vec{F} = F(r)\hat{r}$
    \item Is along the radial direction
    \item Derives from potential: $F(r) = -\frac{dV}{dr}$
\end{itemize}
\end{definition}

\begin{importantbox}
\textbf{Consequences of Central Force:}
\begin{enumerate}
    \item Angular momentum $\vec{L}$ is conserved
    \item Motion is in a plane perpendicular to $\vec{L}$
    \item Areal velocity is constant (Kepler's 2nd law)
    \item Energy is conserved
\end{enumerate}
\end{importantbox}

\subsection{Conservation Laws}

\begin{theorem}[Angular Momentum Conservation]
For central force:
\[
\vec{L} = \vec{r} \times \vec{p} = \text{constant}
\]

Magnitude: $L = mr^2\dot{\theta}$
\end{theorem}

\begin{theorem}[Areal Velocity]
Area swept out per unit time:
\[
\frac{dA}{dt} = \frac{1}{2}r^2\dot{\theta} = \frac{L}{2m} = \text{constant}
\]
\end{theorem}

\subsection{Equation of Motion in Plane}

\begin{formula}[Central Force Equations]
In polar coordinates $(r, \theta)$:

\textbf{Radial equation:}
\[
m\ddot{r} = mr\dot{\theta}^2 - \frac{dV}{dr} = \frac{L^2}{mr^3} - \frac{dV}{dr}
\]

\textbf{Angular equation:}
\[
\frac{d}{dt}(mr^2\dot{\theta}) = 0 \implies L = mr^2\dot{\theta} = \text{const}
\]
\end{formula}

\subsection{Effective Potential}

\begin{definition}
The \textbf{effective potential} is:
\[
V_{\text{eff}}(r) = V(r) + \frac{L^2}{2mr^2}
\]
where the second term is the centrifugal barrier.
\end{definition}

\begin{importantbox}
\textbf{Energy Equation:}
\[
E = \frac{1}{2}m\dot{r}^2 + V_{\text{eff}}(r)
\]

This reduces 2D motion to effective 1D motion in $r$!

\textbf{Analysis:}
\begin{itemize}
    \item $E > V_{\text{eff}}$: Motion possible
    \item $E = V_{\text{eff}}$: Turning point
    \item $\frac{dV_{\text{eff}}}{dr} = 0$ and $E = V_{\text{eff}}$: Circular orbit
\end{itemize}
\end{importantbox}

\subsection{Orbit Equation}

\begin{theorem}[Binet's Equation]
Let $u = 1/r$. The orbit equation is:
\[
\frac{d^2u}{d\theta^2} + u = -\frac{m}{L^2u^2}F(1/u)
\]

For potential $V(r)$:
\[
\frac{d^2u}{d\theta^2} + u = -\frac{m}{L^2u^2}\frac{d}{dr}V(r) = \frac{m}{L^2}\frac{d}{du}V(1/u)
\]
\end{theorem}

\subsection{Kepler Problem (Inverse Square Law)}

\begin{definition}
\textbf{Gravitational force:} $F = -\frac{GMm}{r^2} = -\frac{k}{r^2}$

\textbf{Potential:} $V(r) = -\frac{k}{r}$ where $k = GMm$
\end{definition}

\subsubsection{Orbit Equation for Kepler Problem}

\begin{theorem}
For inverse square law, Binet's equation gives:
\[
\frac{d^2u}{d\theta^2} + u = \frac{mk}{L^2}
\]

General solution:
\[
u = \frac{1}{r} = \frac{mk}{L^2}(1 + e\cos(\theta - \theta_0))
\]

Or, choosing $\theta_0 = 0$:
\[
r = \frac{l}{1 + e\cos\theta}
\]
where $l = \frac{L^2}{mk}$ is the semi-latus rectum and $e$ is eccentricity.
\end{theorem}

\begin{importantbox}
\textbf{Conic Sections Based on Eccentricity:}
\begin{itemize}
    \item $e = 0$: Circle (constant $r$)
    \item $0 < e < 1$: Ellipse (bound orbit)
    \item $e = 1$: Parabola (just unbound)
    \item $e > 1$: Hyperbola (unbound, scattering)
\end{itemize}

Eccentricity in terms of energy:
\[
e = \sqrt{1 + \frac{2EL^2}{mk^2}}
\]
\end{importantbox}

\subsubsection{Kepler's Laws}

\begin{theorem}[Kepler's Three Laws]
\textbf{1. Law of Orbits:} Planets move in ellipses with Sun at one focus.
\[
r = \frac{a(1-e^2)}{1 + e\cos\theta}
\]

\textbf{2. Law of Areas:} Equal areas swept in equal times.
\[
\frac{dA}{dt} = \frac{L}{2m} = \text{const}
\]

\textbf{3. Law of Periods:} $T^2 \propto a^3$
\[
T^2 = \frac{4\pi^2}{GM}a^3
\]
where $a$ is semi-major axis.
\end{theorem}

\begin{example}
\textbf{Derive Kepler's Third Law}

For ellipse: $a = \frac{l}{1-e^2}$, Area $= \pi ab = \pi a^2\sqrt{1-e^2}$

Areal velocity: $\frac{dA}{dt} = \frac{L}{2m}$

Period: $T = \frac{\text{Area}}{dA/dt} = \frac{\pi a^2\sqrt{1-e^2}}{L/2m}$

Using $l = L^2/(mk)$ and $a = l/(1-e^2)$:
\[
T^2 = \frac{4\pi^2 m a^3}{k} = \frac{4\pi^2 a^3}{GM}
\]
\end{example}

\subsection{Energy Analysis}

\begin{formula}[Energy for Different Orbits]
Total energy: $E = \frac{1}{2}m(\dot{r}^2 + r^2\dot{\theta}^2) - \frac{k}{r}$

Using $L = mr^2\dot{\theta}$:
\[
E = \frac{1}{2}m\dot{r}^2 + \frac{L^2}{2mr^2} - \frac{k}{r}
\]

\textbf{For circular orbit} ($\dot{r} = 0$, $r = r_0$):
\[
E = \frac{L^2}{2mr_0^2} - \frac{k}{r_0} = -\frac{k}{2r_0} \quad (\text{using force balance})
\]

\textbf{For elliptical orbit:}
\[
E = -\frac{k}{2a}
\]
where $a$ is semi-major axis.

\textbf{For parabolic orbit:} $E = 0$

\textbf{For hyperbolic orbit:} $E > 0$
\end{formula}

\subsection{Scattering in Central Force}

\begin{definition}
For repulsive central force (e.g., Coulomb scattering), particles follow hyperbolic trajectories.

\textbf{Impact parameter} $b$: Perpendicular distance from force center to incoming trajectory.

\textbf{Scattering angle} $\Theta$: Angle between incoming and outgoing asymptotes.
\end{definition}

\begin{formula}[Rutherford Scattering]
For Coulomb potential $V = k/r$ (repulsive, $k > 0$):
\[
\Theta = \pi - 2\theta_{\min}
\]
where $\theta_{\min}$ is angle at closest approach.

Relation between impact parameter and scattering angle:
\[
b = \frac{k}{2E}\cot\frac{\Theta}{2}
\]

Differential cross-section:
\[
\frac{d\sigma}{d\Omega} = \left(\frac{k}{4E}\right)^2 \frac{1}{\sin^4(\Theta/2)}
\]
\end{formula}

\subsection{Virial Theorem}

\begin{theorem}[Virial Theorem]
For a bound system with potential $V = kr^n$:
\[
\langle T \rangle = \frac{n}{2}\langle V \rangle
\]

For inverse square law ($n = -1$):
\[
\langle T \rangle = -\frac{1}{2}\langle V \rangle \implies E = \langle T \rangle + \langle V \rangle = \frac{1}{2}\langle V \rangle = -\langle T \rangle
\]
\end{theorem}

\subsection{Practice Problems - Central Force Motion}

\begin{enumerate}
    \item Particle in potential $V = -k/r + c/r^2$. Find conditions for circular orbit.
    
    \item Calculate escape velocity from Earth's surface.
    
    \item For elliptical orbit, find ratio of speeds at perihelion and aphelion.
    
    \item Derive scattering angle for hard sphere scattering.
    
    \item Find period of small radial oscillations about circular orbit.
\end{enumerate}

%======================================================================
% CHAPTER 4: RIGID BODY DYNAMICS
%======================================================================
\newpage
\section{Rigid Body Dynamics}

\subsection{Rigid Body Definition and Kinematics}

\begin{definition}
A \textbf{rigid body} is a collection of particles with fixed distances between them:
\[
|r_i - r_j| = \text{constant}
\]

\textbf{Degrees of freedom:} 6 (3 translation + 3 rotation)
\end{definition}

\subsection{Angular Velocity and Rotation}

\begin{definition}
\textbf{Angular velocity vector} $\vec{\omega}$:

Velocity of point $P$ relative to origin:
\[
\vec{v} = \vec{v}_{\text{CM}} + \vec{\omega} \times \vec{r}'
\]
where $\vec{r}'$ is position relative to center of mass.
\end{definition}

\begin{formula}[Rotation Matrices]
Rotation about $z$-axis by angle $\phi$:
\[
R_z(\phi) = \begin{pmatrix}
\cos\phi & -\sin\phi & 0 \\
\sin\phi & \cos\phi & 0 \\
0 & 0 & 1
\end{pmatrix}
\]

Similarly for $R_x(\theta)$ and $R_y(\psi)$.
\end{formula}

\subsection{Euler Angles}

\begin{definition}
\textbf{Euler angles} $(\phi, \theta, \psi)$ specify orientation:
\begin{enumerate}
    \item Rotation by $\phi$ about $z$-axis
    \item Rotation by $\theta$ about new $x$-axis (line of nodes)
    \item Rotation by $\psi$ about new $z$-axis
\end{enumerate}

Total rotation: $R = R_z(\phi)R_x(\theta)R_z(\psi)$
\end{definition}

\begin{importantbox}
\textbf{Angular Velocity in Terms of Euler Angles:}
\begin{align}
\omega_x &= \dot{\phi}\sin\theta\sin\psi + \dot{\theta}\cos\psi \\
\omega_y &= \dot{\phi}\sin\theta\cos\psi - \dot{\theta}\sin\psi \\
\omega_z &= \dot{\phi}\cos\theta + \dot{\psi}
\end{align}
\end{importantbox}

\subsection{Moment of Inertia Tensor}

\begin{definition}
The \textbf{moment of inertia tensor} is:
\[
I_{ij} = \sum_\alpha m_\alpha(r_\alpha^2\delta_{ij} - x_{i\alpha}x_{j\alpha})
\]

In matrix form:
\[
\mathbf{I} = \begin{pmatrix}
I_{xx} & I_{xy} & I_{xz} \\
I_{yx} & I_{yy} & I_{yz} \\
I_{zx} & I_{zy} & I_{zz}
\end{pmatrix}
\]

where:
\begin{align}
I_{xx} &= \sum m(y^2 + z^2) \\
I_{yy} &= \sum m(x^2 + z^2) \\
I_{zz} &= \sum m(x^2 + y^2) \\
I_{xy} &= I_{yx} = -\sum mxy \quad \text{(products of inertia)}
\end{align}
\end{definition}

\begin{importantbox}
\textbf{For Continuous Body:}
\[
I_{xx} = \int (y^2 + z^2)\rho\,dV, \quad I_{xy} = -\int xy\rho\,dV
\]
\end{importantbox}

\subsection{Principal Axes}

\begin{theorem}
Every rigid body has three orthogonal \textbf{principal axes} along which $\mathbf{I}$ is diagonal:
\[
\mathbf{I} = \begin{pmatrix}
I_1 & 0 & 0 \\
0 & I_2 & 0 \\
0 & 0 & I_3
\end{pmatrix}
\]

Principal moments are eigenvalues of inertia tensor.
\end{theorem}

\begin{tipbox}
\textbf{Finding Principal Axes:}
\begin{enumerate}
    \item Form inertia tensor $\mathbf{I}$
    \item Solve $\det(\mathbf{I} - I\lambda) = 0$ for eigenvalues $I_1, I_2, I_3$
    \item Find eigenvectors (principal axes)
\end{enumerate}
\end{tipbox}

\subsection{Parallel Axis Theorem}

\begin{theorem}[Parallel Axis Theorem]
If $I_{\text{CM}}$ is moment of inertia about axis through center of mass, and $I$ is moment about parallel axis at distance $d$:
\[
I = I_{\text{CM}} + Md^2
\]
where $M$ is total mass.
\end{theorem}

\begin{theorem}[Perpendicular Axis Theorem]
For planar body in $xy$-plane:
\[
I_z = I_x + I_y
\]
\end{theorem}

\begin{example}
\textbf{Moments of Inertia for Common Shapes}

\textbf{1. Uniform Rod (length $L$, mass $M$):}
\begin{itemize}
    \item About center, perpendicular: $I = \frac{1}{12}ML^2$
    \item About end, perpendicular: $I = \frac{1}{3}ML^2$
\end{itemize}

\textbf{2. Uniform Disk (radius $R$, mass $M$):}
\begin{itemize}
    \item About center, perpendicular to disk: $I = \frac{1}{2}MR^2$
    \item About diameter: $I = \frac{1}{4}MR^2$
\end{itemize}

\textbf{3. Solid Sphere (radius $R$, mass $M$):}
\begin{itemize}
    \item About diameter: $I = \frac{2}{5}MR^2$
\end{itemize}

\textbf{4. Hollow Sphere (radius $R$, mass $M$):}
\begin{itemize}
    \item About diameter: $I = \frac{2}{3}MR^2$
\end{itemize}

\textbf{5. Solid Cylinder (radius $R$, length $L$, mass $M$):}
\begin{itemize}
    \item About axis: $I = \frac{1}{2}MR^2$
    \item About diameter through center: $I = \frac{1}{12}ML^2 + \frac{1}{4}MR^2$
\end{itemize}
\end{example}

\subsection{Angular Momentum}

\begin{definition}
For rigid body rotating with $\vec{\omega}$:
\[
\vec{L} = \mathbf{I} \cdot \vec{\omega}
\]

In component form: $L_i = \sum_j I_{ij}\omega_j$
\end{definition}

\begin{importantbox}
\textbf{Key Point:} $\vec{L}$ and $\vec{\omega}$ are NOT necessarily parallel unless rotating about principal axis!

In principal axes:
\[
L_1 = I_1\omega_1, \quad L_2 = I_2\omega_2, \quad L_3 = I_3\omega_3
\]
\end{importantbox}

\subsection{Kinetic Energy}

\begin{formula}[Kinetic Energy of Rigid Body]
\[
T = \frac{1}{2}M v_{\text{CM}}^2 + \frac{1}{2}\vec{\omega} \cdot \mathbf{I} \cdot \vec{\omega}
\]

In principal axes:
\[
T_{\text{rot}} = \frac{1}{2}(I_1\omega_1^2 + I_2\omega_2^2 + I_3\omega_3^2)
\]
\end{formula}

\subsection{Euler's Equations}

\begin{theorem}[Euler's Equations of Motion]
In body-fixed principal axes:
\begin{align}
I_1\dot{\omega}_1 + (I_3 - I_2)\omega_2\omega_3 &= N_1 \\
I_2\dot{\omega}_2 + (I_1 - I_3)\omega_3\omega_1 &= N_2 \\
I_3\dot{\omega}_3 + (I_2 - I_1)\omega_1\omega_2 &= N_3
\end{align}
where $N_i$ are components of external torque.
\end{theorem}

\begin{tipbox}
These equations account for the fact that principal axes rotate with the body!
\end{tipbox}

\subsection{Torque-Free Motion}

\begin{theorem}[Torque-Free Symmetric Top]
For symmetric top ($I_1 = I_2 \neq I_3$) with no external torque:

$\omega_3 = $ constant

$\omega_1$ and $\omega_2$ precess:
\[
\omega_1 = A\cos(\Omega t), \quad \omega_2 = A\sin(\Omega t)
\]
where $\Omega = \frac{I_3 - I_1}{I_1}\omega_3$
\end{theorem}

\subsection{Heavy Symmetric Top}

\begin{definition}
Symmetric top with one point fixed (e.g., spinning top under gravity).

Using Euler angles, Lagrangian is:
\[
L = \frac{1}{2}I_1(\dot{\theta}^2 + \dot{\phi}^2\sin^2\theta) + \frac{1}{2}I_3(\dot{\psi} + \dot{\phi}\cos\theta)^2 - Mgl\cos\theta
\]
\end{definition}

\begin{importantbox}
\textbf{Constants of Motion:}
\begin{itemize}
    \item Energy $E$ (time-independent)
    \item $p_\phi = L_z$ (azimuthal angular momentum, $\phi$ cyclic)
    \item $p_\psi = I_3\omega_3$ (spin angular momentum, $\psi$ cyclic)
\end{itemize}

These lead to:
\[
\frac{1}{2}I_1\dot{\theta}^2 + V_{\text{eff}}(\theta) = E
\]
where
\[
V_{\text{eff}}(\theta) = \frac{(p_\phi - p_\psi\cos\theta)^2}{2I_1\sin^2\theta} + Mgl\cos\theta
\]
\end{importantbox}

\subsection{Precession and Nutation}

\begin{definition}
\textbf{Precession:} Slow rotation of symmetry axis around vertical ($\dot{\phi}$)

\textbf{Nutation:} Wobbling of symmetry axis ($\theta$ oscillation)

\textbf{Spin:} Fast rotation about symmetry axis ($\dot{\psi}$)
\end{definition}

\begin{formula}[Steady Precession]
For fast-spinning top ($\omega_3$ large), approximate steady precession rate:
\[
\Omega_p = \dot{\phi} \approx \frac{Mgl}{I_3\omega_3}
\]
\end{formula}

\subsection{Practice Problems - Rigid Body Dynamics}

\begin{enumerate}
    \item Find moment of inertia tensor for rectangular plate.
    
    \item Solid sphere rolls without slipping down incline - find acceleration.
    
    \item Symmetric top spinning - derive Euler's equations and find solutions.
    
    \item Calculate precession frequency of gyroscope.
    
    \item Torque-free asymmetric top - discuss stability of rotation about each principal axis.
\end{enumerate}

%======================================================================
% CHAPTER 5: SMALL OSCILLATIONS
%======================================================================
\newpage
\section{Small Oscillations and Normal Modes}

\subsection{Motivation}

\begin{importantbox}
Small oscillations about equilibrium appear everywhere:
\begin{itemize}
    \item Molecular vibrations
    \item Crystal lattice vibrations (phonons)
    \item Structural engineering
    \item Quantum harmonic oscillators
    \item Perturbations in orbits
\end{itemize}

Method: Expand potential near equilibrium and keep quadratic terms.
\end{importantbox}

\subsection{Single Oscillator Review}

\begin{formula}[Simple Harmonic Oscillator]
Equation: $m\ddot{x} + kx = 0$ or $\ddot{x} + \omega_0^2x = 0$ where $\omega_0 = \sqrt{k/m}$

Solution: $x(t) = A\cos(\omega_0 t + \phi)$

Energy: $E = \frac{1}{2}kA^2 = \frac{1}{2}m\omega_0^2A^2$
\end{formula}

\subsection{Coupled Oscillators - Two Masses}

\begin{example}
\textbf{Two Masses Connected by Springs}

Setup: Masses $m_1$, $m_2$ connected by springs $k_1$, $k$, $k_2$ on frictionless surface.

Displacements: $x_1$, $x_2$ from equilibrium.

Equations of motion:
\begin{align}
m_1\ddot{x}_1 &= -k_1 x_1 + k(x_2 - x_1) \\
m_2\ddot{x}_2 &= -k(x_2 - x_1) - k_2 x_2
\end{align}

Matrix form:
\[
\begin{pmatrix} m_1 & 0 \\ 0 & m_2 \end{pmatrix}
\begin{pmatrix} \ddot{x}_1 \\ \ddot{x}_2 \end{pmatrix}
= -\begin{pmatrix} k_1 + k & -k \\ -k & k + k_2 \end{pmatrix}
\begin{pmatrix} x_1 \\ x_2 \end{pmatrix}
\]

Or: $\mathbf{M}\ddot{\vec{x}} = -\mathbf{K}\vec{x}$
\end{example}

\subsection{General Theory of Small Oscillations}

\begin{theorem}[Small Oscillations About Equilibrium]
For system with generalized coordinates $q_i$, expand potential about equilibrium $q_i^{(0)}$:
\[
V = V_0 + \sum_i \frac{\partial V}{\partial q_i}\bigg|_0 \eta_i + \frac{1}{2}\sum_{ij} \frac{\partial^2 V}{\partial q_i \partial q_j}\bigg|_0 \eta_i\eta_j + \cdots
\]

where $\eta_i = q_i - q_i^{(0)}$ are displacements.

At equilibrium: $\frac{\partial V}{\partial q_i}\bigg|_0 = 0$

Keeping quadratic terms:
\[
V \approx V_0 + \frac{1}{2}\sum_{ij} V_{ij}\eta_i\eta_j
\]
where $V_{ij} = \frac{\partial^2 V}{\partial q_i \partial q_j}\bigg|_0$
\end{theorem}

\begin{formula}[Kinetic Energy]
For small velocities:
\[
T = \frac{1}{2}\sum_{ij} T_{ij}\dot{\eta}_i\dot{\eta}_j
\]
where $T_{ij}$ are mass matrix elements (often $T_{ij} = m_i\delta_{ij}$).
\end{formula}

\subsection{Normal Modes}

\begin{definition}
A \textbf{normal mode} is a motion where all parts oscillate with same frequency and phase.

Try solution: $\eta_i(t) = a_i e^{i\omega t}$

Substituting into $T\ddot{\vec{\eta}} + V\vec{\eta} = 0$:
\[
(-\omega^2 \mathbf{T} + \mathbf{V})\vec{a} = 0
\]

Non-trivial solution requires:
\[
\det(\mathbf{V} - \omega^2\mathbf{T}) = 0
\]

This is the \textbf{secular equation} giving $n$ normal frequencies $\omega_1, \ldots, \omega_n$.
\end{definition}

\begin{importantbox}
\textbf{Procedure to Find Normal Modes:}
\begin{enumerate}
    \item Set up kinetic and potential energy matrices $\mathbf{T}$ and $\mathbf{V}$
    \item Solve secular equation $\det(\mathbf{V} - \omega^2\mathbf{T}) = 0$ for eigenvalues $\omega_k^2$
    \item For each $\omega_k$, solve $(\mathbf{V} - \omega_k^2\mathbf{T})\vec{a}^{(k)} = 0$ for eigenvector $\vec{a}^{(k)}$
    \item General solution is superposition:
    \[
    \eta_i(t) = \sum_k C_k a_i^{(k)}\cos(\omega_k t + \phi_k)
    \]
    \item Determine $C_k$ and $\phi_k$ from initial conditions
\end{enumerate}
\end{importantbox}

\begin{example}
\textbf{Two Identical Masses, Three Identical Springs}

$m_1 = m_2 = m$, $k_1 = k = k_2$

Mass matrix: $\mathbf{M} = m\mathbf{I}$

Stiffness matrix: $\mathbf{K} = k\begin{pmatrix} 2 & -1 \\ -1 & 2 \end{pmatrix}$

Secular equation:
\[
\det\begin{pmatrix} 2k - m\omega^2 & -k \\ -k & 2k - m\omega^2 \end{pmatrix} = 0
\]
\[
(2k - m\omega^2)^2 - k^2 = 0
\]
\[
2k - m\omega^2 = \pm k
\]

Normal frequencies:
\[
\omega_1 = \sqrt{\frac{k}{m}}, \quad \omega_2 = \sqrt{\frac{3k}{m}}
\]

For $\omega_1$: $\vec{a}^{(1)} = \begin{pmatrix} 1 \\ 1 \end{pmatrix}$ (symmetric mode - masses move together)

For $\omega_2$: $\vec{a}^{(2)} = \begin{pmatrix} 1 \\ -1 \end{pmatrix}$ (antisymmetric mode - masses move opposite)

General solution:
\begin{align}
x_1(t) &= A_1\cos(\omega_1 t + \phi_1) + A_2\cos(\omega_2 t + \phi_2) \\
x_2(t) &= A_1\cos(\omega_1 t + \phi_1) - A_2\cos(\omega_2 t + \phi_2)
\end{align}
\end{example}

\subsection{Normal Coordinates}

\begin{definition}
\textbf{Normal coordinates} $Q_k$ are linear combinations of $\eta_i$ such that Lagrangian becomes:
\[
L = \frac{1}{2}\sum_k (\dot{Q}_k^2 - \omega_k^2 Q_k^2)
\]

Each $Q_k$ oscillates independently:
\[
\ddot{Q}_k + \omega_k^2 Q_k = 0
\]
\end{definition}

\begin{formula}[Transformation to Normal Coordinates]
\[
\eta_i = \sum_k a_i^{(k)} Q_k
\]
where $\vec{a}^{(k)}$ are normalized eigenvectors.

In matrix form: $\vec{\eta} = \mathbf{A}\vec{Q}$ where columns of $\mathbf{A}$ are eigenvectors.
\end{formula}

\subsection{Orthogonality of Normal Modes}

\begin{theorem}
Normal modes satisfy orthogonality relations:
\[
\sum_i T_{ij} a_j^{(k)} a_i^{(l)} = M_k\delta_{kl}
\]
\[
\sum_i V_{ij} a_j^{(k)} a_i^{(l)} = M_k\omega_k^2\delta_{kl}
\]

For simple case $T_{ij} = m\delta_{ij}$:
\[
\vec{a}^{(k)} \cdot \vec{a}^{(l)} = \delta_{kl} \quad \text{(after normalization)}
\]
\end{theorem}

\subsection{Energy in Normal Modes}

\begin{formula}
In normal coordinates:
\[
E = \sum_k E_k = \sum_k \frac{1}{2}(\dot{Q}_k^2 + \omega_k^2 Q_k^2)
\]

Each normal mode carries energy independently:
\[
E_k = \frac{1}{2}\omega_k^2 A_k^2
\]
where $Q_k(t) = A_k\cos(\omega_k t + \phi_k)$.
\end{formula}

\subsection{Continuous Systems - String Vibrations}

\begin{example}
\textbf{Vibrating String (Length $L$, Tension $T$, Linear Density $\mu$)}

Wave equation:
\[
\frac{\partial^2 y}{\partial t^2} = v^2 \frac{\partial^2 y}{\partial x^2}
\]
where $v = \sqrt{T/\mu}$ is wave speed.

Boundary conditions: $y(0,t) = y(L,t) = 0$

Normal modes:
\[
y_n(x,t) = A_n\sin\left(\frac{n\pi x}{L}\right)\cos(\omega_n t + \phi_n)
\]

Normal frequencies:
\[
\omega_n = \frac{n\pi v}{L} = \frac{n\pi}{L}\sqrt{\frac{T}{\mu}}, \quad n = 1, 2, 3, \ldots
\]

General solution:
\[
y(x,t) = \sum_{n=1}^{\infty} A_n\sin\left(\frac{n\pi x}{L}\right)\cos(\omega_n t + \phi_n)
\]
\end{example}

\subsection{Damped Oscillations}

\begin{formula}[Damped Harmonic Oscillator]
Equation: $m\ddot{x} + \gamma\dot{x} + kx = 0$ or $\ddot{x} + 2\beta\dot{x} + \omega_0^2 x = 0$

where $\beta = \gamma/(2m)$ is damping coefficient, $\omega_0 = \sqrt{k/m}$.

\textbf{Three cases:}

\textbf{1. Underdamped} ($\beta < \omega_0$): Oscillatory decay
\[
x(t) = Ae^{-\beta t}\cos(\omega t + \phi), \quad \omega = \sqrt{\omega_0^2 - \beta^2}
\]

\textbf{2. Critically damped} ($\beta = \omega_0$): Fastest return without oscillation
\[
x(t) = (A + Bt)e^{-\beta t}
\]

\textbf{3. Overdamped} ($\beta > \omega_0$): Slow exponential decay
\[
x(t) = Ae^{-(\beta - \sqrt{\beta^2 - \omega_0^2})t} + Be^{-(\beta + \sqrt{\beta^2 - \omega_0^2})t}
\]
\end{formula}

\subsection{Forced Oscillations and Resonance}

\begin{formula}[Driven Damped Oscillator]
Equation: $m\ddot{x} + \gamma\dot{x} + kx = F_0\cos(\omega_d t)$

Steady-state solution:
\[
x(t) = A(\omega_d)\cos(\omega_d t - \delta)
\]

Amplitude:
\[
A(\omega_d) = \frac{F_0/m}{\sqrt{(\omega_0^2 - \omega_d^2)^2 + (2\beta\omega_d)^2}}
\]

Phase lag:
\[
\tan\delta = \frac{2\beta\omega_d}{\omega_0^2 - \omega_d^2}
\]

\textbf{Resonance:} Maximum amplitude at $\omega_d = \omega_r = \sqrt{\omega_0^2 - 2\beta^2}$

At resonance:
\[
A_{\max} = \frac{F_0}{2\beta m\omega_0} \approx \frac{F_0}{2\beta m\omega_r} \quad (\text{for small } \beta)
\]

\textbf{Quality factor:}
\[
Q = \frac{\omega_0}{2\beta} = \frac{\text{Energy stored}}{\text{Energy lost per cycle}}
\]
\end{formula}

\subsection{Practice Problems - Small Oscillations}

\begin{enumerate}
    \item Three identical masses connected by identical springs in a line. Find all normal frequencies and modes.
    
    \item Double pendulum with small angles - find normal modes.
    
    \item Molecule vibrations: CO$_2$ (linear) - find normal modes (3 modes).
    
    \item Loaded string: $N$ equal masses on string - find normal frequencies (limit to continuous string).
    
    \item Driven oscillator: find width of resonance peak (full width at half maximum).
\end{enumerate}

%======================================================================
% APPENDIX: EXAM STRATEGY & IMPORTANT FORMULAS
%======================================================================
\newpage
\section{Classical Mechanics - GATE Exam Strategy}

\subsection{Topic-wise Weightage and Strategy}

\begin{tipbox}
\textbf{Expected Question Distribution:}
\begin{itemize}
    \item \textbf{Lagrangian/Hamiltonian:} 2-3 questions (High weightage!)
    \item \textbf{Central Force/Kepler:} 1-2 questions
    \item \textbf{Rigid Body:} 1-2 questions
    \item \textbf{Small Oscillations:} 1-2 questions
\end{itemize}
\end{tipbox}

\subsection{Common Problem Types}

\begin{enumerate}
    \item \textbf{Direct Lagrangian formulation:} Given system, write $L$ and derive equations
    \item \textbf{Constrained motion:} Pendulums, beads on wires, rolling
    \item \textbf{Central force orbits:} Find eccentricity, period, energy
    \item \textbf{Moment of inertia calculations:} Parallel axis, various shapes
    \item \textbf{Normal mode frequencies:} Coupled oscillators, find eigenvalues
\end{enumerate}

\subsection{Essential Formula Quick Reference}

\begin{importantbox}
\textbf{Lagrangian Mechanics:}
\begin{align*}
L &= T - V \\
\frac{d}{dt}\left(\frac{\partial L}{\partial \dot{q}_i}\right) - \frac{\partial L}{\partial q_i} &= 0 \\
p_i &= \frac{\partial L}{\partial \dot{q}_i}
\end{align*}

\textbf{Hamiltonian Mechanics:}
\begin{align*}
H &= \sum p_i\dot{q}_i - L \\
\dot{q}_i &= \frac{\partial H}{\partial p_i}, \quad \dot{p}_i = -\frac{\partial H}{\partial q_i} \\
\{f,g\} &= \sum_i\left(\frac{\partial f}{\partial q_i}\frac{\partial g}{\partial p_i} - \frac{\partial f}{\partial p_i}\frac{\partial g}{\partial q_i}\right)
\end{align*}

\textbf{Central Force:}
\begin{align*}
r &= \frac{l}{1 + e\cos\theta} \quad \text{(orbit equation)} \\
E &= -\frac{k}{2a} \quad \text{(ellipse energy)} \\
T^2 &= \frac{4\pi^2}{GM}a^3 \quad \text{(Kepler's 3rd law)}
\end{align*}

\textbf{Rigid Body:}
\begin{align*}
\vec{L} &= \mathbf{I} \cdot \vec{\omega} \\
T_{\text{rot}} &= \frac{1}{2}\vec{\omega} \cdot \mathbf{I} \cdot \vec{\omega} \\
I &= I_{\text{CM}} + Md^2 \quad \text{(parallel axis)}
\end{align*}

\textbf{Small Oscillations:}
\begin{align*}
\det(\mathbf{V} - \omega^2\mathbf{M}) &= 0 \quad \text{(secular equation)} \\
\omega_{\text{res}} &= \sqrt{\omega_0^2 - 2\beta^2} \quad \text{(resonance)}
\end{align*}
\end{importantbox}

\subsection{Time-Saving Tips}

\begin{enumerate}
    \item Always check for cyclic coordinates first - conserved momenta save time!
    \item For central force, remember $L$ is conserved - reduces to 1D problem
    \item Know moments of inertia by heart for standard shapes
    \item For coupled oscillators, look for symmetry - can simplify eigenvalue problem
    \item In exams, you can often guess normal mode patterns from symmetry
\end{enumerate}

\subsection{Common Mistakes to Avoid}

\begin{enumerate}
    \item Forgetting minus sign in $L = T - V$ (not $T + V$!)
    \item Wrong sign in Hamilton's equations ($\dot{p} = -\partial H/\partial q$)
    \item Confusing angular momentum $L$ with angular velocity $\omega$
    \item Not using parallel axis theorem when needed
    \item Wrong normalization of eigenvectors
    \item Forgetting $r^2\dot{\theta}$ term in polar coordinates kinetic energy
\end{enumerate}

\subsection{Must-Practice Problems}

\begin{enumerate}
    \item Bead on rotating hoop/wire
    \item Double pendulum
    \item Particle in central potential with $V = -k/r + c/r^2$
    \item Atwood machine variations
    \item Rolling cylinders/spheres
    \item Symmetric top under gravity
    \item Three-mass-three-spring system
    \item Forced oscillator at resonance
\end{enumerate}

%======================================================================
% CONCLUSION
%======================================================================
\newpage
\section{Final Words}

Classical Mechanics is the most problem-solving intensive subject. Success requires:

\begin{enumerate}
    \item \textbf{Deep Understanding:} Know why each formulation works
    \item \textbf{Pattern Recognition:} Most problems are variations of standard types
    \item \textbf{Speed:} Practice until you can set up Lagrangian in under 1 minute
    \item \textbf{Accuracy:} Double-check signs and factors of 2!
\end{enumerate}

\begin{importantbox}
\textbf{30-Day Study Plan:}
\begin{itemize}
    \item \textbf{Week 1:} Lagrangian formulation - 20 problems
    \item \textbf{Week 2:} Hamiltonian mechanics - 20 problems
    \item \textbf{Week 3:} Central force + Rigid body - 15 + 15 problems
    \item \textbf{Week 4:} Small oscillations + Mixed practice - 20 problems
\end{itemize}

\textbf{Last 10 Days:} Previous year GATE problems only!
\end{importantbox}

\vspace{1cm}
\begin{center}
\textbf{\Large Master Classical Mechanics = Strong Foundation for All of Physics!}

\vspace{0.5cm}
\textbf{\Large All the Best for GATE 2026!}
\end{center}

\end{document}
